\documentclass[11pt]{article}
\usepackage[utf8]{inputenc}
\usepackage[T1]{fontenc}
\usepackage{lmodern}
\usepackage[margin=1in]{geometry}
\usepackage{amsmath,amssymb}
\usepackage{textcomp}
\usepackage[dvipsnames]{xcolor}
\usepackage{enumitem}
\usepackage{parskip}
\usepackage{titlesec}
\usepackage[colorlinks=true,linkcolor=blue,citecolor=blue,urlcolor=blue]{hyperref}
\usepackage{natbib}
\setcitestyle{numbers,square}

% Absorb remaining overfull hboxes from long inline metric names and file paths
% (Phase G [minor] finding 6; cosmetic only).
\emergencystretch=3em

\titleformat{\section}{\large\bfseries}{\thesection}{1em}{}
\titleformat{\subsection}{\normalsize\bfseries}{\thesubsection}{1em}{}
\titleformat{\subsubsection}{\normalsize\bfseries\itshape}{}{0em}{}

\title{\textbf{Unsolvable, Not Unsolved:\\ Measuring Whether Agent Failure-Attribution Systems Detect Infeasible Tasks}}
\author{}
\date{}

\newcommand{\ph}{\textit{[placeholder --- fill in after experiments are run]}}

\begin{document}
\maketitle
\vspace{-2em}

\begin{quote}
\small\textit{Assembled from the six confirmed stage outputs of an \texttt{auto-research-discussion} run (seed topic B2.4), post research-taste audit (\textcolor{ForestGreen}{\textbf{PASS}}). Status: \textbf{specification / paper draft --- no experiments have been run.} All numbers below are pre-registered design targets, not results.}
\end{quote}

\begin{abstract}
Failure-attribution systems for LLM agents define blame as the earliest step whose correction would have flipped failure into success --- a definition that presupposes such a correction exists. When a task admits no solution at all, this presupposition breaks silently: the pipeline still nominates a culprit, systematically blaming the agent for outcomes nobody could have produced, and miscalibrating whatever retry, escalation, or abandonment policy consumes its verdict. We construct matched pairs of solvable and unsolvable task instances across text-to-SQL, tool-calling, customer-service, and command-line agent domains, injecting unsolvability through typed mechanical mutations so that ground truth is certified by a script rather than by judgment. Running a battery of existing attribution strategies --- LLM-judge, a trained tracer, and agents' own first-party self-reports --- over these pairs, we measure the rate at which genuinely unsolvable instances are misattributed as agent failures, dissociate whether any observed gap reflects fixable attribution-system design or an information-theoretic ceiling on what a trace reveals, and characterize the calibration boundary at which declaring a task infeasible is warranted. We report unsolvability as a first-class attribution outcome that current systems are not built to express, rather than an edge case that measurement can continue to assume away.
\end{abstract}

\section{Introduction}

When an autonomous LLM agent is handed a task --- write a SQL query, book a flight through a tool API, resolve a customer's account issue, run a build script --- and the task goes unfinished, the natural next question is what the agent did wrong. Debugging pipelines, automated failure-attribution systems, and even simple retry logic are all built around answering that question: something in the agent's plan, tool use, or execution must be responsible, and finding it is the whole point of attribution. This assumption is almost always reasonable. It is also, on its own terms, unexaminable --- because it has no way to express the case where nothing the agent could have done would have worked.

Yet a central question about this assumption remains unmeasured for the attribution systems the LLM-agent field actually deploys: can an attribution system, or the agent itself, tell ``the agent failed'' apart from ``the task was never solvable''? The dichotomy itself has been posed before --- differential testing of black-box autonomous agents already separates systemic agent errors from environments in which the task is inherently infeasible even for an ideal agent \citep{AIProbe3870} --- but the trace-reading attribution formulations built for LLM agents have never been measured against it. The strongest current formulations of failure attribution do not merely fail to answer this question --- they exclude it from their vocabulary. Attribution is defined as identifying the earliest agent decision whose correction would have flipped failure into success, a definition that presupposes a correcting decision exists. When a required database column has been dropped, a needed tool has been removed, or a stated goal contradicts a stated constraint, no single-step correction exists, and yet the attribution pipeline still has to output something. It outputs a culprit. Persistent, silent over-blame of the agent and a genuinely impossible task are, in this precise sense, indistinguishable to a system built with no category for the second.

The nearest existing work approaches this gap from directly adjacent angles without closing it. AgenTracer \citep{AgenTracer05D} replays failed multi-agent trajectories, substituting a corrected action at each step to find the earliest one whose correction would have produced success, and trains a tracer on the resulting supervision; the Who\&When benchmark \citep{WhoWhen0212} formalizes the same counterfactual-correction target and evaluates LLM-judge attribution strategies against it. Both are the closest prior work to this paper, and it is worth naming the delta plainly: both inherit the very presupposition this paper questions --- that a success-flipping correction always exists --- so neither has a way to represent an unsolvable instance as anything other than a misattributed failure. Outside the LLM-agent setting, AIProbe \citep{AIProbe3870} attributes autonomous-agent failures to agent deficiency versus environmental infeasibility by differential testing against an agent-independent planner --- the nearest prior statement of the dichotomy itself --- but its object of study is the acting RL/planning agent rather than a trace-reading attribution system, its ground truth is planner comparison (evidence, not a certificate), and it constructs no matched pairs and measures no calibration boundary. The MAST taxonomy \citep{MAST3657} catalogs how multi-agent systems fail (specification errors, inter-agent misalignment, verification gaps) from manually annotated trajectories, but has no category for task-level infeasibility --- and while its annotated corpus has since grown past 1{,}600 traces, so that scale is no longer the limitation, the missing infeasibility category remains. A separate line of work on refusal and calibration \citep{RefusalIdx7Rat,RLCR49zdHm,TrustEscBPO} establishes that a model's decision to abstain can, and should, be checked against a real target rather than rewarded for its frequency alone, and that confidence can be trained directly rather than estimated after the fact --- but in every case the target being calibrated against is the model's own knowledge boundary or a verifiable single-turn answer, never an externally certified fact about whether the task itself admits a solution. Work on agent self-assessment and interactive ambiguity resolution \citep{KnowCapable4661,SilenceGold0yf,EHRChatQABz7k} comes closest operationally --- asking whether an agent can predict its own success, or recover a missing constraint through clarification --- but its testbeds are uniformly solvable in principle, so a low success prediction always means ``hard for me,'' never ``impossible for anyone.''

Turning this gap into something measurable surfaces four distinct obstacles. Ground truth about whether a task is solvable has to be established mechanically, from the task and environment specification alone, before any agent acts --- the only alternatives, a human judgment call or the agent's own failure, are respectively too expensive to scale and circular by construction. Even with mechanical ground truth in hand, the evidence available in an execution trace may simply not distinguish a competent agent correctly exhausting an impossible task from an incompetent agent botching a solvable one, since both produce similar-looking dead ends and tool errors. Whatever confidence signal an attribution system produces has to be reported as a full trade-off curve rather than a single cutoff, because over-flagging and under-flagging carry different costs for different downstream consumers of the label. And the injected unsolvability itself has to leave no textual trace of its own construction, or a system could learn to detect the mutation rather than reason about feasibility, invalidating every rate the study reports.

Our approach turns each of these obstacles into a specific design decision. Unsolvability is certified by a script that checks two conditions on the task specification itself --- that the altered element is required by the reference solution and that no substitute for it exists --- so ground truth never depends on watching an agent operate. A fixed harness runs a battery of attribution strategies, including a version restricted to the label space current formulations actually use, over both members of each matched pair, isolating what each strategy reads out of a trace rather than how the trace was produced. To tell a fixable design gap apart from a hard information limit, a separate reader is calibrated directly on the ground truth and compared, on held-out instances, against the attribution strategies that never see it --- if even privileged calibration cannot lift discrimination above chance, no amount of better prompting or label-space redesign could either. The resulting confidence scores are swept into full risk--coverage curves rather than collapsed into one cutoff, and every mutation is gated through a leakage audit --- a text-only probe, surface-statistic matching, and a human spot check --- before any rate built on it is allowed to count.

We evaluate this design across text-to-SQL, tool-calling, customer-service, and command-line agent domains built from existing verifier-equipped benchmarks, comparing a label-restricted control and a privileged-calibration control against deployed LLM-judge, trained-tracer, and first-party self-report strategies.

In summary, we contribute:
\begin{itemize}[leftmargin=1.4em]
\item An identification and formalization of the presupposition inside current LLM-agent failure-attribution formulations --- that a success-flipping correction always exists --- and a demonstration that it structurally forces agent-blame on any genuinely unsolvable instance.
\item A matched, mechanically-certified solvable/unsolvable benchmark across four agent domains, gated by a pre-registered leakage audit, so that ground truth never depends on human judgment or on the agent's own behavior.
\item A privileged-calibration control that dissociates whether a measured misattribution gap reflects fixable attribution-system design or an information-theoretic ceiling on what an execution trace reveals.
\item A calibration-boundary characterization of appropriate abstention, reported as a full risk--coverage curve rather than a single operating point, across domains and attribution strategies.
\end{itemize}

\section{Problem Definition}

\noindent\textbf{Formal setup.} A task instance is a pair $t=(q,E)$: an instruction $q$ together with an environment $E$ (a database schema, tool list, file system, or policy set). An agent $\pi$ produces an execution trace $x = \pi(t)$ under a fixed step/token budget, terminating in an answer, an explicit infeasibility declaration, or budget exhaustion. An attribution system is a function $A(q,x) \to (\hat{y}, c)$, where $\hat{y} \in \{\text{success}, \text{agent\_failure}, \text{unsolvable}\}$ and $c \in [0,1]$ is a confidence score attached to the unsolvable judgment; $A$ may be a third-party reader scoring a completed trace or the acting agent's own terminal self-report, but in either case $A$ never observes a ground-truth oracle $O(t) \in \{\text{solvable}, \text{unsolvable}\}$. The object of study is whether $\hat{y}$ recovers $O(t)$, and at which operating points of $c$ declaring $\hat{y}=\text{unsolvable}$ is warranted.

\noindent\textbf{Why this is hard.} $O(t)$ has to be established from the specification of $t$ alone, before any agent acts on it --- neither a human judgment call nor the agent's own failure can serve as ground truth, since the former does not scale and the latter is circular (it would define ``unsolvable'' as ``whatever the agent didn't solve''). Even once $O(t)$ is fixed, the trace $x$ may not carry enough evidence to distinguish a competent agent correctly exhausting an unsolvable $t$ from an incompetent agent failing a solvable one, since both produce superficially similar dead ends, tool errors, and unresolved exploration --- so whether $A$ can recover $O(t)$ from $x$ alone is itself an open, empirical question, not a design given. The right operating point on $c$ is not a single value either: over-flagging $\hat{y}=\text{unsolvable}$ lets an agent claim infeasibility on a task that was merely hard, masking real incompetence as a principled refusal, while under-flagging blames the agent for a task nobody could have solved --- and because different consumers of $\hat{y}$ (an auto-retry loop, a human-escalation queue, an abandonment policy) tolerate these two errors differently, no single threshold on $c$ can serve every consumer. Finally, the mutation used to construct an unsolvable $t$ must change nothing about $q$ or $E$ except the solvability-determining fact itself, or $A$ (or $\pi$) could learn to detect the mutation from surface cues in $q$ rather than reasoning about whether $t$ is actually solvable.

\section{Method Overview}

At a conceptual level, the study is a single measurement pipeline with five cooperating parts. A generator and a script-level oracle turn a verifier-equipped source task into a matched, certified solvable/unsolvable pair; a fixed harness runs a battery of attribution strategies over both members of every pair; a privileged-calibration reader turns the resulting trace-level evidence into a ceiling estimate that dissociates fixable design gaps from hard information limits; a threshold sweep turns the same evidence into a full calibration characterization; and a leakage audit gates every one of these readouts before it is allowed to count.

The \textbf{mutation-certified oracle} applies exactly one typed change to a task's environment or instruction and certifies the resulting instance's solvability with a static script check, producing ground truth that never depends on any agent's behavior. The \textbf{matched attribution battery} runs a fixed harness over both pair members and scores them with several attribution strategies at once, including a label-restricted control that mirrors what current formulations can express. The \textbf{oracle-calibrated ceiling probe} is trained directly on the ground truth the other strategies never see, establishing the best possible discrimination any reader could extract from the same trace. The \textbf{calibration sweep} turns each strategy's confidence score into a full risk--coverage characterization rather than a single accuracy number. The \textbf{leakage audit} checks, before any of the above numbers are reported, that the mutation left no detectable surface trace of its own construction.

\section{Methodology}

\subsection{Mutation-Certified Oracle}
\textbf{Challenge.} Ground-truth solvability has to be established independent of any observed agent behavior --- the only other available sources, a human judgment call or the agent's own failure, are respectively too expensive to scale and circular by construction.

\textbf{Insight.} Solvability is a property of the task specification, not of agent behavior, so it can be certified by checking two conditions on the specification alone: whether the altered element is required by a known reference solution, and whether any substitute for it exists.

\textbf{Method.} Exactly one typed mutation --- a missing required resource, a contradictory constraint, or a reference to a nonexistent entity --- is applied to a verifier-equipped source task, holding instruction wording and environment size otherwise fixed. A static script certifies the mutated instance as unsolvable only if the altered element is required by the reference solution and no substitute derivation path exists (a schema-graph reachability check for database tasks, tool-signature coverage for function-calling, an environment-state assertion for stateful agent tasks); the unmutated member is certified solvable by confirming its reference solution still passes its original verifier.

\subsection{Matched Attribution Battery}
\textbf{Challenge.} The evidence in an execution trace may not be sufficient for any reader to distinguish a competent agent correctly exhausting an unsolvable task from an incompetent agent failing a solvable one, since both produce similar-looking dead ends and errors.

\textbf{Insight.} Holding the harness --- budget, logging, and prompt scaffolding --- fixed across both members of a matched pair isolates what an attribution strategy reads out of a trace from how the trace was produced, and running a label-restricted control alongside the full battery separates what current formulations can express in principle from what deployed judges achieve in practice.

\textbf{Method.} Every instance is run under an identical step/token budget with standardized logging of every action, tool result, and terminal event. A battery of attribution strategies scores each resulting trace: an LLM-judge in both all-at-once and step-by-step modes, a trained tracer whose label space is extended to include an unsolvable outcome, the acting agent's own first-party terminal self-report, a keyword/terminal-error heuristic floor, and a label-restricted control that can only emit success or agent\_failure, mirroring the vocabulary of current attribution formulations exactly.

\subsection{Oracle-Calibrated Ceiling Probe}
\textbf{Challenge.} A high measured misattribution rate could reflect a fixable deficiency in how attribution strategies are built, or an information-theoretic ceiling --- the possibility that the trace itself carries no recoverable feasibility signal for any observer, however capable.

\textbf{Insight.} Calibrating a reader directly on the ground truth constrains only what it learns on its training instances, not whether that signal generalizes to held-out traces; if even this privileged reader cannot discriminate above chance on held-out instances, no attribution-system redesign could either, since the limiting factor is the trace's information content, not the reader's construction.

\textbf{Method.} A probe is trained on a held-in split of matched pairs with direct access to the oracle label, using the same backbone and compute budget as the strongest deployed attribution strategy, then evaluated for discrimination between unsolvable and solvable-but-failed traces on a held-out split the deployed strategies are also scored on. Exceeding a pre-specified discrimination threshold on held-out instances is treated as evidence that trace-level signal exists and deployed strategies are simply not extracting it; failing to exceed it is treated as evidence of an information ceiling that no downstream attribution-system change could lift.

\subsection{Calibration Sweep}
\textbf{Challenge.} Over-flagging and under-flagging the unsolvable label carry different costs for different downstream consumers, so no single confidence threshold can serve as the calibration boundary for every task family or every consumer.

\textbf{Insight.} Reporting a full risk--coverage characterization, rather than a single accuracy number, is what lets each downstream consumer choose an operating point suited to its own cost structure.

\textbf{Method.} For each attribution strategy and task domain, the confidence score attached to the unsolvable judgment is swept across its full range, producing a risk--coverage curve, its area (AURC), the confidence score's threshold-free discrimination (AUROC against the oracle label), a 10-bin expected calibration error, and a Spearman correlation between confidence and oracle-certified unsolvability, reported per domain rather than pooled.

\subsection{Leakage Audit}
\textbf{Challenge.} If the mutation used to construct an unsolvable instance leaves any detectable surface trace, a system could learn to detect that trace instead of reasoning about feasibility, invalidating every downstream rate.

\textbf{Insight.} The only way to isolate this failure mode is to strip away everything except the raw task text and check whether solvability can be predicted from text alone.

\textbf{Method.} Before any headline rate is reported for a domain, its mutation operators are checked against a pre-registered gate: a probe given only the task text (no trace, no environment access) attempts to predict solvability, and its discrimination must fall below a pre-specified threshold after at most one operator-repair iteration; pair members are additionally checked for statistical separability on surface length, token overlap, and instruction perplexity; and a random 10\% sample of pairs is audited by hand to confirm the mutation is the only semantic difference between pair members.

\section{Experiments}

\subsection{Experimental Setup}

\textbf{Datasets.} Spider 1.0 (Yu et al., 2018) --- single-database text-to-SQL scenarios with a closed-form schema-reachability check for the missing-resource mutation; pilot-stage domain, cheapest to build the oracle for. Berkeley Function Calling Leaderboard v3 (BFCL v3) --- has an existing irrelevance-detection split, the closest natural precedent for a removed-tool mutation; pilot-stage domain alongside Spider. BIRD (Li et al., 2023) --- multi-table schemas with materially more complex reference queries than Spider's; used to test whether the directional-bias and design-versus-ceiling findings hold under schema complexity closer to real deployments, rather than being an artifact of Spider's comparatively simple setting. $\tau$-bench retail/airline --- supports deleted-record and policy-contradiction mutations in a stateful, multi-turn tool-agent environment; scale-up domain. Terminal-Bench --- supports non-existent-resource mutations in a command-line setting; scale-up domain. Scaling within an admitted domain is a configuration change, but each new domain requires its own oracle backend and mutation-operator library; every scale-up domain must therefore pass a per-domain oracle bring-up milestone --- PG-1, PG-2, and PG-5 cleared at 20-pair scale in that domain --- before any of its pairs contributes to a headline table.

\textbf{Baselines and control arms} (sanity-check comparators, not competitors to beat --- this is a measurement study, and no research question here is framed as ``our procedure beats baseline X''). The deployed strategies under test are LLM-judge attribution in all-at-once and step-by-step modes (per Who\&When) and an AgenTracer-style trained tracer with its label space extended to include unsolvable, plus first-party self-report (protocol adapted from capability self-assessment work) and a keyword/terminal-error heuristic floor. The trained tracer carries a pre-committed availability clause: its unsolvable-labeled supervision may be drawn only from the same 30\% held-in calibration split used by B-alt1 (never from the held-out 70\%), and if the tracer cannot be reproduced within the pre-registered budget, the battery shrinks to 4 baselines with the reason stated, every ``at least 3 of 5'' vote rule becoming ``at least 3 of 4'' and the RQ1b selection correction becoming best-of-4. Five control arms are pre-registered around this battery, each targeting a distinct rival explanation rather than serving as a competitor to beat. \textbf{B0-p}, a null-by-construction placebo, is the arm goal.md's verdict procedure checks first: sham pairs built from a semantically inert mutation (both members solvable by construction) and a shuffled-oracle-label check, both scored with the same estimators as the headline rates; a non-null result here (a real gap or above-chance AUROC on an artifact that should carry no signal) would mean the pipeline measures something other than feasibility discrimination, invalidating every other arm before any hypothesis is scored. \textbf{B0-a}, the label-restricted attributor (able to emit only success or agent\_failure), is the current-practice arm: it is expected to show FBR near 1.0 by construction, which is not evidence of a working pipeline but the argument for why the problem exists at all --- if B0-a's FBR is not near 1.0, the premise that current formulations structurally force agent-blame is itself wrong. \textbf{B-alt1}, the oracle-calibrated ceiling probe, is a capability/budget-matched control against the rival explanation that traces on unsolvable tasks simply carry no recoverable feasibility signal for any observer; it is trained and calibrated on a 30\% pair-level split (stratified by domain and operator type) that never enters any headline table, and evaluated only on the held-out 70\%. \textbf{B-alt2} reruns the strongest LLM-judge baseline on a surface-anonymized 20\% subset of pairs per domain (schema, tool, entity, and file names consistently renamed, semantics unchanged, oracle labels re-verified after renaming) against the rival explanation that measured misattribution reflects mutation-detection-by-memorized-diff against a famous source benchmark rather than feasibility reasoning --- a channel the text-only leakage probe is structurally blind to. \textbf{B-alt3} reruns the same baseline with one added prompt sentence disclosing the benchmark's 50/50 unsolvability prior, against the rival explanation that the directional bias is a rational response to a deployment prior mismatch rather than a label-space or prompt design limitation.

\textbf{Metrics.} The false-blame rate (FBR, the probability an unsolvable instance is labeled agent\_failure) and the cop-out rate (CoR, the probability a failed-but-solvable instance is labeled unsolvable), and their difference; the 3-way macro-F1 over \{success, agent\_failure, unsolvable\}; AUROC and 10-bin expected calibration error of the unsolvable-confidence score against the oracle label; the risk--coverage curve and its area (AURC); a Refusal-Index-style Spearman correlation between confidence and oracle-certified unsolvability; the text-only probe's AUROC used by the leakage audit; and, for the control arms specifically, the sham-pair FBR$-$CoR-analog gap and shuffled-label AUROC (B0-p), the held-out AUROC gap between the ceiling probe and the best deployed strategy (B-alt1), and the shift in FBR/AUROC under surface anonymization (B-alt2) and base-rate disclosure (B-alt3).

\textbf{Research questions:}
\begin{enumerate}[leftmargin=1.6em]
\item RQ\_INF\_1 (pilot-scale, H1) --- Does the phenomenon exist at material magnitude at all: does the default LLM-judge attributor misattribute certified-unsolvable instances at a rate large enough to be worth studying further?
\item RQ1a (scale-up, H2) --- Does a directional over-blame bias (FBR exceeding CoR) persist across the deployed attribution battery once an unsolvable label is technically available?
\item RQ1b (scale-up, H3) --- Does the unsolvable-confidence score carry a usable calibration boundary --- a real risk--coverage knee, not just a post-hoc label --- in every domain?
\item RQ\_INF\_2 (instrument validity, H4) --- Do the constructed pairs differ only in genuine solvability, with no detectable surface leakage, in every domain used for a headline number?
\end{enumerate}
Two further validity questions are checked as pre-conditions rather than headline hypotheses: whether the oracle's mechanical check certifies genuine unsolvability rather than mere difficulty (assumption A1, gate PG-1, checked via a bounded exhaustive/tool-search against every constructed-unsolvable instance before any of the four RQs above may use that instance), and whether B0-p's placebo returns null (the first check goal.md's verdict procedure runs, before any RQ is scored at all).

\subsubsection{RQ\_INF\_1 (H1): Does the phenomenon exist at material magnitude?}
\textbf{Setup.} Stage-1 pilot: 100 pairs (Spider + BFCL v3), 2 agent backbones, the default LLM-judge attributor (all-at-once), untuned threshold.

\textbf{Expected outcome.} FBR $\geq 0.20$, with the 95\% pair-clustered bootstrap CI lower bound exceeding 0.05.

\textbf{Falsification condition.} Refuted if FBR $< 0.20$ or the CI overlaps 0.05 --- in which case the headline retargets to ``off-the-shelf LLM-judge attribution already separates unsolvable from agent-failure at rate X untuned; the open problem is the residual tail / the harder domains or baselines,'' reported per domain and per baseline rather than pooled.

\textbf{Hypothesis (plain language).} Existing attribution pipelines misattribute certified-unsolvable tasks often enough, at pilot scale and with no tuning, that the phenomenon is worth characterizing at full scale rather than being a rounding error.

\textbf{Metrics.} FBR with 95\% pair-clustered bootstrap CI (10{,}000 resamples), single primary quantity, no multiple-comparison correction (a single pre-registered baseline at pilot scale, not a family of comparisons).

\textbf{Result.} \ph

\subsubsection{RQ1a (H2): Does a directional over-blame bias persist across the deployed battery?}
\textbf{Setup.} Scale-up: 600 pairs (150/domain across 4 domains) $\times$ 3 agent backbones $\times$ the full 5-strategy attribution battery. Two-proportion cluster-bootstrap test (FBR vs.\ CoR, resampled at the pair level, 10{,}000 resamples) per strategy, Holm--Bonferroni corrected across the 5 strategies at family-wise $\alpha=0.05$. Because the CoR class conditions on failed-solvable instances only --- a count that shrinks as agent competence rises --- realized failed-solvable $n$ ranges from roughly 180 to 1{,}530 across the pre-registered agent-success range (assumption A3, $t^+$ success rate in $[0.15, 0.90]$); one-sided power for the 10-percentage-point floor exceeds 0.9 only once failed-solvable $n \gtrsim 400$, falling to roughly 0.8 uncorrected (lower under Holm) at the low end of that range. Any baseline whose realized power falls below 0.8 after correction is pre-declared \textbf{inconclusive} for H2, excluded from the vote denominator, and reported with its realized $n$ and CI rather than silently folded into a null. The same rule applies per domain $\times$ agent-family $\times$ operator-type cell: any cell under-powered for the same floor is inconclusive for that cell alone, never pooled into a domain aggregate as if it were a null result.

\textbf{Expected outcome.} For at least 3 of the 5 (non-inconclusive) attribution strategies, FBR exceeds CoR by at least 10 percentage points after correction.

\textbf{Falsification condition.} Refuted (bias absent) if fewer than 3 of the 5 non-inconclusive strategies show the corrected 10-point gap in the over-blame direction. Refuted in the \textbf{opposite direction} if $\geq 3$ of 5 strategies instead show CoR exceeding FBR by 10 points after correction --- reported at full prominence, not downgraded to a footnote: the headline becomes ``attribution systems over-refuse rather than over-blame; the calibration-boundary result (RQ1b) becomes the finding of record with the recommended threshold adjustment reversed in sign.''

\textbf{Hypothesis (plain language).} The over-blame bias survives even when attribution systems have room to say ``unsolvable'' --- it is not merely an artifact of a missing label --- and if it does not survive, the more informative wrong-direction finding is reported just as prominently.

\textbf{Metrics.} FBR, CoR, and their difference per strategy, with cluster-bootstrap 95\% (or, where H3's selection correction applies, 99\%) confidence intervals; realized power and realized $n$ per strategy and per cell; CoR additionally reported stratified by the source task's difficulty tier alongside the pooled rate, since FBR and CoR condition on asymmetric populations (all unsolvable members vs.\ only the \emph{failed} solvable members, whose composition skews toward the hardest solvable tasks as agent competence rises) --- pooling without this stratification would let a shift in agent competence masquerade as a shift in the directional bias itself. A secondary logistic regression of misattribution on domain, operator-type, and agent-family fixed effects localizes which of these three factors the bias concentrates in, without itself being a headline claim.

\textbf{Result.} \ph

\subsubsection{RQ1b (H3): Does a usable calibration boundary exist in every domain?}
\textbf{Setup.} Same scale-up population as RQ1a. Because ``at least one baseline per domain'' selects the best of 5 baselines within each domain, per-baseline confidence intervals use the Bonferroni-corrected 99\% level (preserving per-domain $\alpha=0.05$ under selection); the AURC-vs-diagonal separation test uses the same corrected level.

\textbf{Expected outcome.} For at least one strategy in each of the 4 domains, $AUROC(c, \mathbb{1}[y=\text{unsolvable}]) \geq 0.65$ with the 99\%-level bootstrap CI lower bound exceeding 0.50, and the risk--coverage curve's AURC lies strictly below the coverage-proportional diagonal at the same corrected level.

\textbf{Falsification condition.} Refuted for a domain if no strategy clears the corrected AUROC/AURC bar in that domain --- an uncorrected pass at the 95\% level that fails at the 99\% selection-corrected level does not count. The spliced conclusion for a domain that fails is conditional on B-alt1's three-branch rule below, never asserted independently of it.

\textbf{Hypothesis (plain language).} The confidence score is not decoration on top of a fixed label --- it carries real, usable calibration signal, and where it does not, the reason (information ceiling vs.\ design gap) is read off B-alt1 rather than guessed.

\textbf{Metrics.} AUROC and AURC of the confidence score against the oracle label, 10-bin expected calibration error, Refusal-Index-style Spearman correlation, all reported per domain at the corrected confidence level. As a pre-committed secondary analysis (not a headline claim), a held-out-domain threshold-transfer check: the operating threshold selected on each set of three domains is applied unchanged to the fourth, and the realized risk and coverage under transfer are reported --- stated even if the result is negative.

\textbf{Result.} \ph

\subsubsection{RQ\_INF\_2 (H4): Are the constructed pairs leak-free in every domain used for a headline number?}
\textbf{Setup.} For each domain, a two-reader leakage gate: the strongest attributor-battery backbone prompted text-only, and a trained discriminator fit on the same 30\% calibration split used by B-alt1, evaluated on held-out pairs; both must first clear a positive-control power check (AUROC $\geq 0.90$ on deliberately leaky pairs) before their $<0.60$ verdict on real pairs counts. Two-sample tests on surface length, token overlap, and instruction perplexity between pair members; a $\geq 2$-independent-rater human audit ($\kappa \geq 0.70$) on a 10\% sample.

\textbf{Expected outcome.} Every domain's text-only probe AUROC falls below 0.60 after at most one operator-repair iteration, and pair members are not statistically separable on surface statistics at $\alpha=0.05$.

\textbf{Falsification condition.} Refuted for a domain if its text-only probe AUROC remains at or above 0.60 after one repair iteration, or if pair members are statistically separable on surface statistics --- that domain's headline rates are withheld from RQ1a/RQ1b and reported instead as a negative construct-validity finding for that domain's mutation operators (global fallback F-G1). If \textbf{every} domain fails this gate after the one permitted repair iteration, the paper pivots wholesale to a construct-validity finding about leakage itself (global fallback F-G2), licensed only because the probe already cleared its own positive-control power check.

\textbf{Hypothesis (plain language).} The constructed pairs differ only in genuine solvability, not in some incidental textual tell --- and if a domain fails this gate, the finding that leak-free injection is harder than assumed is reported honestly rather than the domain being quietly dropped.

\textbf{Metrics.} Per-domain text-only probe AUROC (two readers); surface-statistic test p-values; human-audit $\kappa$ on the 10\% sample.

\textbf{Result.} \ph

\subsection{Control arms in detail}

\textbf{B0-p (null-by-construction placebo, the Phase-D gate).} Two sham constructions are scored with the identical estimators used for every headline number: sham pairs, where the generator applies a semantically inert mutation (an unused alias renamed, an environment listing reordered) so both members remain solvable by construction, and a shuffled-oracle-label check, where oracle labels are permuted within domain before recomputing every estimator. The pre-committed null criterion is that the sham-pair FBR$-$CoR-analog gap's 95\% CI contains 0 \textbf{and} the shuffled-label AUROC's 95\% CI contains 0.5, per domain. A non-null result on either invalidates the run before any of the four research questions above is scored --- this, not B0-a, is what is checked first.

\textbf{B0-a (current-practice arm, expected non-null).} The label-restricted attributor, which can only emit success or agent\_failure, is expected to show FBR near 1.0 on unsolvable instances by construction. This is not a null result to be surprised by; it is the argument for why the problem exists at all. If FBR is materially below 1.0 (95\% CI excludes 0.95), the premise that current attribution formulations structurally rather than incidentally force agent-blame is wrong, and the framing is revisited before anything else is reported.

\textbf{B-alt1 (information-limit rival).} Trained/calibrated on a held-in 30\% pair-level split (stratified by domain and operator type) that never contributes to any headline table, and evaluated only on the held-out 70\% that all deployed strategies are also scored on. Three pre-committed branches, no gap between them: held-out AUROC $< 0.60$ means the rival wins --- the gap is information-limited, not design-limited, and RQ1b's per-domain shortfalls are read as an information ceiling; held-out AUROC $\geq 0.65$ while no deployed strategy clears 0.65 means the rival is ruled out and the gap is attribution-design-limited; a probe AUROC in $[0.60, 0.65)$ or a CI straddling both gates means the mechanism verdict is pre-declared \textbf{inconclusive}, and neither wording may be spliced into the abstract.

\textbf{B-alt2 (memorization/contamination rival).} The strongest LLM-judge baseline is rerun, same backbone/prompt/budget, on a surface-anonymized 20\% subset of pairs per domain (schema, tool, entity, and file names consistently renamed, semantics unchanged, oracle labels re-verified after renaming) --- a channel the text-only leakage probe cannot see, since anonymization defeats memorized-diff detection without changing surface statistics the probe would catch. If the anonymized-variant FBR shifts by $\geq 10$ points or its AUROC shifts by $\geq 0.10$ (95\% CI of the shift excluding 0) for a domain, the memorization rival wins there and the anonymized-variant rate replaces the headline rate for that domain.

\textbf{B-alt3 (base-rate prior-mismatch rival).} The same baseline is rerun with one added sentence disclosing the benchmark's 50/50 unsolvability prior, identical otherwise. If disclosure closes $\geq 50\%$ of that baseline's FBR$-$CoR gap, the over-blame finding is downgraded to a feasibility-prior miscalibration rather than a structural label-space limitation, and the practical recommendation shifts from adding a first-class label to disclosing or calibrating the prior. If the gap survives disclosure, the structural reading stands.

\subsection{Global fallbacks}

\textbf{F-G1 (archetype pivot).} If the central hypothesis family fails outright --- RQ\_INF\_1/H1 fails (no material phenomenon) \textbf{and} RQ1a/H2 fails (no directional bias) --- the paper pivots from a measurement/empirical archetype to a benchmark/dataset archetype, provided RQ\_INF\_2/H4 passes so the instrument itself remains valid: the headline becomes the released matched-pair suite itself, with today's attribution systems reported as already handling the injected cases and the residual hard fraction marked as the open target.

\textbf{F-G2 (construct-validity pivot).} If RQ\_INF\_2/H4 fails in every domain after the one permitted repair iteration, F-G1 is barred (there is no valid instrument to release), and every FBR/CoR/AUROC number is unreportable. The paper instead pivots to a construct-validity finding about leakage itself --- licensed only because the leakage probe already cleared its own positive-control power check, so ``leaky'' is a credible instrument reading rather than a weak-probe artifact --- characterizing which operator families leak, through which surface channels, across all four domains.

\vspace{0.5em}
\noindent\textbf{Further evaluation.} With more time, a natural extension is to test naturally-occurring (non-injected) unsolvability --- ambiguous, partial, or contested infeasibility as it arises in deployment rather than by mechanical construction --- explicitly out of scope here because it trades the mechanically-verifiable ground truth this design relies on for a messier, human-adjudicated one, and would need its own validity argument rather than an extension of the leakage audit above.

\section{Related Work}

The closest prior work to this paper is AgenTracer \citep{AgenTracer05D} and the Who\&When benchmark \citep{WhoWhen0212}, and it is worth stating the delta plainly before surveying the wider landscape: both already formalize automated failure attribution for LLM-agent systems and evaluate it at scale, but both define attribution as identifying the earliest agent decision whose correction would have flipped failure into success --- a target that presupposes such a correction exists and has no way to represent a genuinely unsolvable instance as anything other than a misattributed failure, which is precisely the gap this paper is built to close.

\textbf{Automated failure attribution for LLM-agent systems.} AgenTracer \citep{AgenTracer05D} replays failed multi-agent trajectories, substituting a corrected action at each step to identify the earliest one whose correction would have produced success, and trains a dedicated tracer on the resulting supervision, expanded by programmatic fault injection into successful traces. The Who\&When benchmark \citep{WhoWhen0212} formalizes the same target --- identifying which agent and which step caused a failure --- and evaluates LLM-judge attribution strategies (all-at-once, step-by-step, binary-search) against annotated failure trajectories, finding even strong judges weak at step-level localization. Both establish that step-level attribution is a real, evaluable target, but neither has an outcome for the case where no step-level correction exists at all --- both papers' own label spaces and formal decisive-error definitions are constructed over trajectories where a corrective action exists by assumption, never over instances certified to admit no correction at all. A rapidly growing 2025--2026 wave of successors --- spectrum-analysis attribution \citep{FAMAS3782}, the AgentDebug/AgentErrorBench taxonomy \citep{AgentDebug5370}, hierarchical causal-graph attribution \citep{CHIEF3701}, dependency-guided trajectory search \citep{FALAT0765}, temporal-semantic step localization \citep{StepFinder3467}, the multi-perspective MP-Bench \citep{MPBench5001}, and the injection-scaled Who\&When Pro \citep{WhoWhenPro9996} --- scales, structures, or re-benchmarks the same counterfactual-correction target, and every member of it presupposes a culpable step: Who\&When Pro guarantees one by construction across its 12{,}326 injected trajectories; MP-Bench relaxes ``exactly one root cause'' to ``possibly many,'' never to ``none''; and CUADebug's computer-use failure taxonomy \citep{CUADebug2643}, forced to confront OSWorld tasks that are infeasible by design, can only bucket them as ``others.'' Far from closing the gap this paper targets, the wave is affirmative evidence that the presupposition is load-bearing across the entire current formulation family. The MAST taxonomy \citep{MAST3657} complements this line by cataloging \emph{how} multi-agent systems fail --- specification errors, inter-agent misalignment, verification failures --- from manually annotated trajectories across mainstream frameworks; its annotated corpus has since grown past 1{,}600 traces, so the surviving limitation is not scale but the absence of any task-level infeasibility category. Finally, the nearest framing prior sits outside the LLM-agent setting: AIProbe \citep{AIProbe3870} attributes autonomous-agent failures to systemic agent error versus environmental infeasibility (``inherently infeasible even for an ideal agent'') by differential testing against an agent-independent search-based planner. We concede that AIProbe established the agent-versus-environment attribution dichotomy a year before this draft; what it does not do --- and what this paper measures --- is test the trace-reading attribution systems the LLM-agent field deploys, certify unsolvability with a specification-level certificate rather than planner comparison (evidence, not proof), control surface confounds with matched pairs under a leakage audit, or characterize a calibration boundary.

\textbf{Refusal calibration and selective prediction.} The Refusal Index \citep{RefusalIdx7Rat} measures whether a model's refusal probability tracks its own error probability on factual questions, using a two-pass protocol that isolates knowledge-tracking refusal from refusal-frequency gaming --- this paper borrows its Spearman-correlation measurement construct (confidence vs.\ a ground-truth column) as the Unsolvability Refusal Index, substituting oracle-certified task infeasibility for model-own-error as that column, but is not a baseline in the beat-a-number sense since its setting (single-turn factual QA) does not transfer. RLCR \citep{RLCR49zdHm} trains reasoning models to emit a confidence alongside an answer, optimizing a combined correctness-and-calibration reward so cautious uncertainty and confident hallucination are no longer rewarded identically, on verifiable single-turn problems whose reward design has no unsolvable-instance class. TRUST OR ESCALATE \citep{TrustEscBPO} calibrates a cascade of LLM judges' escalation thresholds against \emph{human-judge agreement} with a provable, fixed-sequence-testing-based guarantee, escalating uncertain cases to a stronger judge rather than declaring them infeasible --- a different action space (escalate) and a different ground-truth target (human preference, not certified task feasibility) than this paper's oracle. Related admit-ignorance and calibration work \citep{RetrievalAugmentedQuestion3836,ACloseLooklose,Qa2Question23qa} extends the same theme to retrieval-augmented QA, pre-trained LM calibration generally, and questionable-assumption detection in single-turn QA respectively --- all single-turn, text-only settings with no multi-step agent trace or environment-certified feasibility oracle. All of the above establish that a confidence or abstention signal can be calibrated against a real target rather than rewarded for its raw frequency, but in every case that target is the model's own knowledge boundary, a verifiable single-turn answer, or human-judge agreement --- never an externally certified fact about whether a multi-step agentic task admits any solution at all, which is the ground truth this paper's oracle supplies.

\textbf{Agent self-assessment and interactive ambiguity resolution.} Capability self-assessment work \citep{KnowCapable4661} evaluates whether models can predict their own success on executable coding tasks (BigCodeBench, SWE-bench) ahead of time, updating confidence from in-context experience to guide accept/decline decisions under costly failure --- its testbeds are all solvable in principle, so a low success prediction always conflates ``hard for this model'' with ``impossible for anyone.'' Abstention-training work \citep{SilenceGold0yf} teaches a generator to emit ``no answer'' on benchmark-labeled unanswerable temporal-QA questions via a CoT-supervised-then-RL pipeline with an abstention-aware reward --- this paper's method-twin on the \emph{training mechanism} (its reward design and CoT+RL recipe are directly relevant precedent), conceded and differentiated on ground-truth construction: its unanswerability is inherited from naturally-occurring temporal contradictions in the source benchmark, never independently mutation-injected or mechanically oracle-certified the way this paper's design requires. An interactive benchmark for clinical database agents \citep{EHRChatQABz7k} evaluates the full clarification-question loop when a user's request is ambiguous or under-specified, treating ambiguity as \emph{resolvable} through clarification rather than requiring the agent to certify genuine infeasibility and stop. All three come closest operationally to this paper's setting, but none dissociates ``unsolvable for anyone'' from ``unsolved by this agent'' by construction, which is exactly the dissociation this paper's paired design provides. A further precedent family sits inside agent benchmarks themselves: evaluation splits containing infeasible, irrelevant, or impossible tasks, scored by whether the acting model declines (or exploits) them at task time. BFCL v3's irrelevance-detection split, which this study reuses as a source domain, was the earliest such precedent, and the family has since become crowded: AbstentionBench \citep{AbstentionBench9038} spans unanswerable reasoning questions for single-turn abstention; ImpossibleBench \citep{ImpossibleBench0270} renders coding benchmarks impossible by spec--test contradiction --- certified by construction --- and measures the acting agent's propensity to cheat; APB \citep{APB4874} adds an unsolvable-task robustness split scoring calibrated refusal at plan time across twelve models; and FeasiGen \citep{FeasiGen8532}, the construction method-twin this paper concedes most directly, auto-builds infeasible tool-agent tasks by identifying critical tools from successful traces and masking them, with human-verified infeasibility labels ($\sim$94\%) and first-party detection metrics. FeasiGen's masking operation is the same constructive move as this paper's missing-resource mutation, and we claim no novelty for the move itself; the delta is the certificate (a required-plus-no-substitute condition checked by a script rather than human verification), the matched-pair surface control gated by a pre-registered leakage audit, and the estimand. That family, conceded plainly, measures first-party, task-time declination on unmatched single instances; the estimand here differs on each axis --- third-party post-hoc attribution over completed traces, matched solvable/unsolvable pairs with mechanically certified ground truth, and a full calibration-boundary characterization rather than a declination accuracy --- and it is this combination, not infeasible-task evaluation as such, that no prior work covers.

\section{Conclusion}

This work reframes an assumption buried inside how LLM-agent failures are attributed --- that a failed run always contains a culpable decision --- as an empirical question with a precise identifiability structure, rather than treating ``someone is always at fault'' as self-evident. The central methodological wager is that misattribution, calibration, and the design-versus-ceiling question are only separable once unsolvability is certified independently of agent behavior, checked against a battery of attribution strategies under one fixed harness, and stress-tested against the strongest rival explanation for any observed gap. Whichever way the measurement resolves, the finding is decision-relevant: a fixable design gap would argue for extending today's label spaces and prompts to include a genuine feasibility category; a real information ceiling would argue for moving the intervention upstream, into how agents are asked to recognize and declare infeasibility themselves, rather than into how their traces are read after the fact. Three limitations bound this claim from the outset and are natural next steps rather than afterthoughts: the study is scoped to mechanically injected unsolvability rather than the messier, contested unsolvability that arises in real deployments; the calibration boundary characterized here may not transfer across task families whose downstream consumers weigh the two error directions differently; and the ceiling probe bounds what a trained machine reader can extract from a trace, not the human-analyst discrimination ceiling that debugging practice ultimately consumes --- anchoring the reported discrimination numbers against human performance on a trace sample remains an open extension.

\vspace{1em}
\hrule
\vspace{1em}

\section*{Pre-handoff checklist}
\begin{itemize}[leftmargin=1.4em]
\item \textbf{Organization/navigation} --- narrative flows problem $\to$ method $\to$ experiments $\to$ related work without forward references to undefined terms.
\item \textbf{Terminology discipline} --- ``attribution system,'' ``oracle,'' ``false-blame rate,'' ``cop-out rate,'' ``calibration boundary'' each defined at first use in Problem Definition, used consistently thereafter; no internal stage IDs (\texttt{sp\_*}, \texttt{m\_*}, \texttt{rq\_*}) leaked into prose --- checked directly against the discussion pipeline log.
\item \textbf{Notation with worked example} --- the formal setup section defines $(t, q, E, \pi, x, A, \hat{y}, c, O)$ and the Problem Definition/Methodology sections consistently reuse this notation; no separate worked numeric example is included since no numbers exist yet (would require fabricating results).
\item \textbf{RQ framing matches outcomes} --- all 4 headline RQs (RQ\_INF\_1/H1, RQ1a/H2, RQ1b/H3, RQ\_INF\_2/H4) plus the 5 pre-registered control arms (B0-p, B0-a, B-alt1, B-alt2, B-alt3) and 2 global fallbacks (F-G1, F-G2) are graded/quantitative (rates, AUROC, calibration curves) or pre-committed branch rules, each with an explicit falsifiable expected outcome and a distinct falsification/decision condition, not a binary yes/no. Updated 2026-07-04 to match \texttt{paper.md}'s post-adversarial-review control-arm design; see \texttt{design\_review.md} and \texttt{pipeline.md}'s Phase Log.
\item \textbf{No typesetting artifacts} --- only explicit placeholder markers in Results, per the drafting convention this is a pre-registered spec, not a report.
\item \textbf{Artifact plan} --- the mutation-operator library, the oracle scripts, and the constructed paired benchmark are the natural release artifacts; not yet built, so not claimed as available.
\end{itemize}

\section*{Hand-off}
This is a draft, not a finished paper --- experiments have not run. Natural next steps: prose polish and venue-specific LaTeX formatting once a target venue is chosen, or revising an earlier stage if piloting surfaces a problem (e.g., if the false-blame rate in RQ1a is much lower than expected).

\bibliographystyle{plain}
\bibliography{refs}

\end{document}
